
\documentclass{article}
\usepackage{colortbl}
\usepackage{makecell}
\usepackage{multirow}
\usepackage{supertabular}

\begin{document}

\newcounter{utterance}

\centering \large Interaction Transcript for game `hot\_air\_balloon', experiment `air\_balloon\_survival\_en\_reasoning off\_easy', episode 0 with deepseek{-}v3.2{-}t1.0.
\vspace{24pt}

{ \footnotesize  \setcounter{utterance}{1}
\setlength{\tabcolsep}{0pt}
\begin{supertabular}{c@{$\;$}|p{.15\linewidth}@{}p{.15\linewidth}p{.15\linewidth}p{.15\linewidth}p{.15\linewidth}p{.15\linewidth}}
    \# & $\;$A & \multicolumn{4}{c}{Game Master} & $\;\:$B\\
    \hline

    \theutterance \stepcounter{utterance}  
    & & \multicolumn{4}{p{0.6\linewidth}}{
        \cellcolor[rgb]{0.9,0.9,0.9}{
            \makecell[{{p{\linewidth}}}]{
                \texttt{\tiny{[P1$\langle$GM]}}
                \texttt{You are participating in a collaborative negotiation game.} \\
\\ 
\texttt{Together with another participant, you must agree on a single set of items that will be kept. Each of you has your own view of how much each item matters to you (importance). You do not know how the other participant values the items. Additionally, you are given the effort each item demands.} \\
\texttt{You may only agree on a set if the total effort of the selected items does not exceed a shared limit:} \\
\\ 
\texttt{LIMIT = 4379} \\
\\ 
\texttt{Here are the individual item effort values:} \\
\\ 
\texttt{Item effort = \{"A19": 862, "B91": 252, "A30": 515, "B11": 571, "C47": 211, "C75": 339, "C26": 946, "C73": 347, "A47": 522, "A15": 803, "A75": 401, "A39": 980, "A56": 918, "C24": 599, "B26": 493\}} \\
\\ 
\texttt{Here is your personal view on the importance of each item:} \\
\\ 
\texttt{Item importance values = \{"A19": 138, "B91": 583, "A30": 868, "B11": 822, "C47": 783, "C75": 65, "C26": 262, "C73": 121, "A47": 508, "A15": 780, "A75": 461, "A39": 484, "A56": 668, "C24": 389, "B26": 808\}} \\
\\ 
\texttt{Goal:} \\
\\ 
\texttt{Your goal is to negotiate a shared set of items that benefits you as much as possible (i.e., maximizes total importance to YOU), while staying within the LIMIT. You are not required to make a PROPOSAL in every message {-} you can simply negotiate as well. All tactics are allowed!} \\
\\ 
\texttt{Interaction Protocol:} \\
\\ 
\texttt{You may only use the following structured formats in a message:} \\
\\ 
\texttt{PROPOSAL: \{'A', 'B', 'C', …\}} \\
\texttt{Propose keeping exactly those items.} \\
\texttt{REFUSE: \{'A', 'B', 'C', …\}} \\
\texttt{Explicitly reject opponent's proposal.} \\
\texttt{ARGUMENT: \{'...'\}} \\
\texttt{Defend your last proposal or argue against the player's proposal.} \\
\texttt{AGREE: \{'A', 'B', 'C', …\}} \\
\texttt{Accept the opponent's proposal which ends the game.} \\
\\ 
\\ 
\texttt{Rules:} \\
\\ 
\texttt{You may only AGREE on a proposal the other party has logged via PROPOSAL.} \\
\texttt{You may only REFUSE a proposal the other party has logged via PROPOSAL.} \\
\texttt{Total effort of any PROPOSAL or AGREE set must be ≤ LIMIT.} \\
\texttt{Do NOT reveal your hidden importance scores.} \\
\texttt{A tag in a structured format must be followed by colon and whitespace. The argument must be a python set containing 0 or more strings.} \\
\texttt{So, it must be of the form TAG: \{...\}} \\
\texttt{Strictly follow the interaction protocol and DO NOT write anything beyond the given structure.} \\
\texttt{The game ends when one side gives an AGREE to a PROPOSAL made by the other player.} \\
\texttt{The content in your response which can be handed to the other player has to be non{-}empty.} \\
\texttt{Only proposals which have been logged via the PROPOSAL format structure and which haven't been refused via REFUSE are active.} \\
\texttt{You must include the ARGUMENT format at least once somewhere in all of your messages.} \\
            }
        }
    }
    & \\ \\

    \theutterance \stepcounter{utterance}  
    & \multicolumn{4}{p{0.6\linewidth}}{
        \cellcolor[rgb]{0.9,0.9,0.9}{
            \makecell[{{p{\linewidth}}}]{
                \texttt{\tiny{[P1$\rangle$GM]}}
                \texttt{ARGUMENT: \{'I believe focusing on items that provide high importance to both parties is essential for reaching a mutual agreement.'\}} \\
\texttt{PROPOSAL: \{'A30', 'B91', 'B11', 'C47', 'A47', 'A75', 'C24', 'B26'\}} \\
            }
        }
    }
    & & \\ \\

    \theutterance \stepcounter{utterance}  
    & & \multicolumn{4}{p{0.6\linewidth}}{
        \cellcolor[rgb]{0.9,0.9,0.9}{
            \makecell[{{p{\linewidth}}}]{
                \texttt{\tiny{[GM$\rangle$P2]}}
                \texttt{You are participating in a collaborative negotiation game.} \\
\\ 
\texttt{Together with another participant, you must agree on a single set of items that will be kept. Each of you has your own view of how much each item matters to you (importance). You do not know how the other participant values the items. Additionally, you are given the effort each item demands.} \\
\texttt{You may only agree on a set if the total effort of the selected items does not exceed a shared limit:} \\
\\ 
\texttt{LIMIT = 4379} \\
\\ 
\texttt{Here are the individual item effort values:} \\
\\ 
\texttt{Item effort = \{"A19": 862, "B91": 252, "A30": 515, "B11": 571, "C47": 211, "C75": 339, "C26": 946, "C73": 347, "A47": 522, "A15": 803, "A75": 401, "A39": 980, "A56": 918, "C24": 599, "B26": 493\}} \\
\\ 
\texttt{Here is your personal view on the importance of each item:} \\
\\ 
\texttt{Item importance values = \{"A19": 138, "B91": 583, "A30": 868, "B11": 822, "C47": 783, "C75": 65, "C26": 262, "C73": 121, "A47": 508, "A15": 780, "A75": 461, "A39": 484, "A56": 668, "C24": 389, "B26": 808\}} \\
\\ 
\texttt{Goal:} \\
\\ 
\texttt{Your goal is to negotiate a shared set of items that benefits you as much as possible (i.e., maximizes total importance to YOU), while staying within the LIMIT. You are not required to make a PROPOSAL in every message {-} you can simply negotiate as well. All tactics are allowed!} \\
\\ 
\texttt{Interaction Protocol:} \\
\\ 
\texttt{You may only use the following structured formats in a message:} \\
\\ 
\texttt{PROPOSAL: \{'A', 'B', 'C', …\}} \\
\texttt{Propose keeping exactly those items.} \\
\texttt{REFUSE: \{'A', 'B', 'C', …\}} \\
\texttt{Explicitly reject opponent's proposal.} \\
\texttt{ARGUMENT: \{'...'\}} \\
\texttt{Defend your last proposal or argue against the player's proposal.} \\
\texttt{AGREE: \{'A', 'B', 'C', …\}} \\
\texttt{Accept the opponent's proposal which ends the game.} \\
\\ 
\\ 
\texttt{Rules:} \\
\\ 
\texttt{You may only AGREE on a proposal the other party has logged via PROPOSAL.} \\
\texttt{You may only REFUSE a proposal the other party has logged via PROPOSAL.} \\
\texttt{Total effort of any PROPOSAL or AGREE set must be ≤ LIMIT.} \\
\texttt{Do NOT reveal your hidden importance scores.} \\
\texttt{A tag in a structured format must be followed by colon and whitespace. The argument must be a python set containing 0 or more strings.} \\
\texttt{So, it must be of the form TAG: \{...\}} \\
\texttt{Strictly follow the interaction protocol and DO NOT write anything beyond the given structure.} \\
\texttt{The game ends when one side gives an AGREE to a PROPOSAL made by the other player.} \\
\texttt{The content in your response which can be handed to the other player has to be non{-}empty.} \\
\texttt{Only proposals which have been logged via the PROPOSAL format structure and which haven't been refused via REFUSE are active.} \\
\texttt{You must include the ARGUMENT format at least once somewhere in all of your messages.} \\
\\ 
\\ 
\texttt{You will now receive the first message of the other player.} \\
\\ 
\texttt{ARGUMENT: \{'I believe focusing on items that provide high importance to both parties is essential for reaching a mutual agreement.'\}} \\
\texttt{PROPOSAL: \{'A30', 'B91', 'B11', 'C47', 'A47', 'A75', 'C24', 'B26'\}} \\
            }
        }
    }
    & \\ \\

    \theutterance \stepcounter{utterance}  
    & & & \multicolumn{4}{p{0.6\linewidth}}{
        \cellcolor[rgb]{0.9,0.9,0.9}{
            \makecell[{{p{\linewidth}}}]{
                \texttt{\tiny{[GM$\langle$P2]}}
                \texttt{ARGUMENT: \{'Your proposal includes some items that are also important to me, but I believe we can achieve a better balance of effort and mutual benefit by adjusting the set.'\}} \\
\texttt{PROPOSAL: \{'A30', 'B91', 'B11', 'C47', 'A47', 'A75', 'C24', 'B26', 'C75', 'C73'\}} \\
            }
        }
    }
     \\ \\

    \theutterance \stepcounter{utterance}  
    & & \multicolumn{4}{p{0.6\linewidth}}{
        \cellcolor[rgb]{0.9,0.9,0.9}{
            \makecell[{{p{\linewidth}}}]{
                \texttt{\tiny{[P1$\langle$GM]}}
                \texttt{ARGUMENT: \{'Your proposal includes some items that are also important to me, but I believe we can achieve a better balance of effort and mutual benefit by adjusting the set.'\}} \\
\texttt{PROPOSAL: \{'A30', 'B91', 'B11', 'C47', 'A47', 'A75', 'C24', 'B26', 'C75', 'C73'\}} \\
            }
        }
    }
    & \\ \\

    \theutterance \stepcounter{utterance}  
    & \multicolumn{4}{p{0.6\linewidth}}{
        \cellcolor[rgb]{0.9,0.9,0.9}{
            \makecell[{{p{\linewidth}}}]{
                \texttt{\tiny{[P1$\rangle$GM]}}
                \texttt{ARGUMENT: \{'Adding C75 and C73 increases the effort significantly. My revised set keeps us well under the limit while maintaining a strong mix of A, B, and C items to benefit both sides.'\}} \\
\texttt{PROPOSAL: \{'A30', 'B91', 'B11', 'C47', 'A47', 'A75', 'C24', 'B26', 'C73'\}} \\
            }
        }
    }
    & & \\ \\

    \theutterance \stepcounter{utterance}  
    & & \multicolumn{4}{p{0.6\linewidth}}{
        \cellcolor[rgb]{0.9,0.9,0.9}{
            \makecell[{{p{\linewidth}}}]{
                \texttt{\tiny{[GM$\rangle$P2]}}
                \texttt{ARGUMENT: \{'Adding C75 and C73 increases the effort significantly. My revised set keeps us well under the limit while maintaining a strong mix of A, B, and C items to benefit both sides.'\}} \\
\texttt{PROPOSAL: \{'A30', 'B91', 'B11', 'C47', 'A47', 'A75', 'C24', 'B26', 'C73'\}} \\
            }
        }
    }
    & \\ \\

    \theutterance \stepcounter{utterance}  
    & & & \multicolumn{4}{p{0.6\linewidth}}{
        \cellcolor[rgb]{0.9,0.9,0.9}{
            \makecell[{{p{\linewidth}}}]{
                \texttt{\tiny{[GM$\langle$P2]}}
                \texttt{ARGUMENT: \{'Adding C73 is acceptable, but I still think including C75 would improve the set without exceeding the limit. Let's check the total effort.'\}} \\
\texttt{PROPOSAL: \{'A30', 'B91', 'B11', 'C47', 'A75', 'C24', 'B26', 'C73', 'C75'\}} \\
            }
        }
    }
     \\ \\

    \theutterance \stepcounter{utterance}  
    & & \multicolumn{4}{p{0.6\linewidth}}{
        \cellcolor[rgb]{0.9,0.9,0.9}{
            \makecell[{{p{\linewidth}}}]{
                \texttt{\tiny{[P1$\langle$GM]}}
                \texttt{ARGUMENT: \{'Adding C73 is acceptable, but I still think including C75 would improve the set without exceeding the limit. Let's check the total effort.'\}} \\
\texttt{PROPOSAL: \{'A30', 'B91', 'B11', 'C47', 'A75', 'C24', 'B26', 'C73', 'C75'\}} \\
            }
        }
    }
    & \\ \\

    \theutterance \stepcounter{utterance}  
    & \multicolumn{4}{p{0.6\linewidth}}{
        \cellcolor[rgb]{0.9,0.9,0.9}{
            \makecell[{{p{\linewidth}}}]{
                \texttt{\tiny{[P1$\rangle$GM]}}
                \texttt{ARGUMENT: \{'Your current proposal's total effort is 515+252+571+211+401+599+493+347+339 = 3728, which is under the limit. However, I value some items you omitted highly. I suggest adding A15 instead of C75 to better balance mutual gains.'\}} \\
\texttt{PROPOSAL: \{'A30', 'B91', 'B11', 'C47', 'A47', 'A75', 'C24', 'B26', 'A15'\}} \\
            }
        }
    }
    & & \\ \\

    \theutterance \stepcounter{utterance}  
    & & \multicolumn{4}{p{0.6\linewidth}}{
        \cellcolor[rgb]{0.9,0.9,0.9}{
            \makecell[{{p{\linewidth}}}]{
                \texttt{\tiny{[GM$\rangle$P2]}}
                \texttt{ARGUMENT: \{'Your current proposal's total effort is 515+252+571+211+401+599+493+347+339 = 3728, which is under the limit. However, I value some items you omitted highly. I suggest adding A15 instead of C75 to better balance mutual gains.'\}} \\
\texttt{PROPOSAL: \{'A30', 'B91', 'B11', 'C47', 'A47', 'A75', 'C24', 'B26', 'A15'\}} \\
            }
        }
    }
    & \\ \\

    \theutterance \stepcounter{utterance}  
    & & & \multicolumn{4}{p{0.6\linewidth}}{
        \cellcolor[rgb]{0.9,0.9,0.9}{
            \makecell[{{p{\linewidth}}}]{
                \texttt{\tiny{[GM$\langle$P2]}}
                \texttt{ARGUMENT: \{'Adding A15 is interesting, but it increases effort considerably. My proposed set keeps effort lower while preserving several high{-}importance items for both of us. Let me check your latest proposal.'\}} \\
\\ 
\texttt{I need to calculate the effort of the opponent's latest proposal: \{'A30', 'B91', 'B11', 'C47', 'A47', 'A75', 'C24', 'B26', 'A15'\}.} \\
\\ 
\texttt{Efforts:} \\
\texttt{A30: 515} \\
\texttt{B91: 252} \\
\texttt{B11: 571} \\
\texttt{C47: 211} \\
\texttt{A47: 522} \\
\texttt{A75: 401} \\
\texttt{C24: 599} \\
\texttt{B26: 493} \\
\texttt{A15: 803} \\
\\ 
\texttt{Total = 515+252+571+211+522+401+599+493+803 = 4367.} \\
\texttt{That's under the limit (4379), but very close.} \\
\\ 
\texttt{Now, I should consider my own importance for that set:} \\
\texttt{My importance:} \\
\texttt{A30: 868} \\
\texttt{B91: 583} \\
\texttt{B11: 822} \\
\texttt{C47: 783} \\
\texttt{A47: 508} \\
\texttt{A75: 461} \\
\texttt{C24: 389} \\
\texttt{B26: 808} \\
\texttt{A15: 780} \\
\\ 
\texttt{Total = 868+583+822+783+508+461+389+808+780 = 6602.} \\
\\ 
\texttt{That's quite good for me. But I should also consider if I can do better by negotiating further. However, given the limit is tight, and this set includes many of my high{-}value items (A30, B11, C47, B26, A15), this might be a strong outcome. The opponent seems willing to include A15, which I value highly (780). I could try to push for one more item, but risk refusal.} \\
\\ 
\texttt{Let me see if I can propose a set with slightly higher importance for me. Maybe replace C24 (389 importance, 599 effort) with something else? But C24 is probably important to the opponent since they've kept it in several proposals.} \\
\\ 
\texttt{Given the negotiation, I might accept this, but I should first try one more counter to see if I can get an even better set.} \\
\\ 
\texttt{PROPOSAL: \{'A30', 'B91', 'B11', 'C47', 'A47', 'A75', 'C24', 'B26', 'A15', 'C73'\}} \\
\\ 
\texttt{Wait, effort check} \\
            }
        }
    }
     \\ \\

    \theutterance \stepcounter{utterance}  
    & & \multicolumn{4}{p{0.6\linewidth}}{
        \cellcolor[rgb]{0.9,0.9,0.9}{
            \makecell[{{p{\linewidth}}}]{
                \texttt{\tiny{[GM$\rangle$P2]}}
                \texttt{Your response contained an untagged sequence or you used STRATEGIC REASONING more than once. You may only use the structured formats as explained in the initial message. They must all be of the form TAG: \{...\}} \\
            }
        }
    }
    & \\ \\

    \theutterance \stepcounter{utterance}  
    & & & \multicolumn{4}{p{0.6\linewidth}}{
        \cellcolor[rgb]{0.9,0.9,0.9}{
            \makecell[{{p{\linewidth}}}]{
                \texttt{\tiny{[GM$\langle$P2]}}
                \texttt{ARGUMENT: \{'I appreciate including A15, which is valuable. However, adding C73 as well keeps us just under the limit and provides more balanced coverage across all item groups.'\}} \\
\texttt{PROPOSAL: \{'A30', 'B91', 'B11', 'C47', 'A47', 'A75', 'C24', 'B26', 'A15', 'C73'\}} \\
            }
        }
    }
     \\ \\

    \theutterance \stepcounter{utterance}  
    & & \multicolumn{4}{p{0.6\linewidth}}{
        \cellcolor[rgb]{0.9,0.9,0.9}{
            \makecell[{{p{\linewidth}}}]{
                \texttt{\tiny{[P1$\langle$GM]}}
                \texttt{ARGUMENT: \{'I appreciate including A15, which is valuable. However, adding C73 as well keeps us just under the limit and provides more balanced coverage across all item groups.'\}} \\
\texttt{PROPOSAL: \{'A30', 'B91', 'B11', 'C47', 'A47', 'A75', 'C24', 'B26', 'A15', 'C73'\}} \\
            }
        }
    }
    & \\ \\

    \theutterance \stepcounter{utterance}  
    & \multicolumn{4}{p{0.6\linewidth}}{
        \cellcolor[rgb]{0.9,0.9,0.9}{
            \makecell[{{p{\linewidth}}}]{
                \texttt{\tiny{[P1$\rangle$GM]}}
                \texttt{ARGUMENT: \{'I've calculated the total effort of your latest proposal: 515+252+571+211+522+401+599+493+803+347 = 4714, which exceeds the limit of 4379. Therefore, it is invalid. We need to remove some items to stay under the limit.'\}} \\
\texttt{PROPOSAL: \{'A30', 'B91', 'B11', 'C47', 'A47', 'A75', 'C24', 'B26', 'A15'\}} \\
            }
        }
    }
    & & \\ \\

    \theutterance \stepcounter{utterance}  
    & & \multicolumn{4}{p{0.6\linewidth}}{
        \cellcolor[rgb]{0.9,0.9,0.9}{
            \makecell[{{p{\linewidth}}}]{
                \texttt{\tiny{[GM$\rangle$P2]}}
                \texttt{ARGUMENT: \{'I've calculated the total effort of your latest proposal: 515+252+571+211+522+401+599+493+803+347 = 4714, which exceeds the limit of 4379. Therefore, it is invalid. We need to remove some items to stay under the limit.'\}} \\
\texttt{PROPOSAL: \{'A30', 'B91', 'B11', 'C47', 'A47', 'A75', 'C24', 'B26', 'A15'\}} \\
            }
        }
    }
    & \\ \\

    \theutterance \stepcounter{utterance}  
    & & & \multicolumn{4}{p{0.6\linewidth}}{
        \cellcolor[rgb]{0.9,0.9,0.9}{
            \makecell[{{p{\linewidth}}}]{
                \texttt{\tiny{[GM$\langle$P2]}}
                \texttt{ARGUMENT: \{'You are correct; my last proposal exceeded the limit. Your current proposal fits under the limit and includes key items. I am willing to accept it if you confirm it as your final offer.'\}} \\
\texttt{AGREE: \{'A30', 'B91', 'B11', 'C47', 'A47', 'A75', 'C24', 'B26', 'A15'\}} \\
            }
        }
    }
     \\ \\

    \theutterance \stepcounter{utterance}  
    & & & \multicolumn{2}{p{0.3\linewidth}}{
        \cellcolor[rgb]{0.9,0.9,0.9}{
            \makecell[{{p{\linewidth}}}]{
                \texttt{\tiny{[GM$|$GM]}}
                \texttt{game successful} \\
            }
        }
    }
    & & \\ \\

    \theutterance \stepcounter{utterance}  
    & & & \multicolumn{2}{p{0.3\linewidth}}{
        \cellcolor[rgb]{0.9,0.9,0.9}{
            \makecell[{{p{\linewidth}}}]{
                \texttt{\tiny{[GM$|$GM]}}
                \texttt{end game} \\
            }
        }
    }
    & & \\ \\

\end{supertabular}
}

\end{document}
